\section{Related Work}
\label{sec:related-work}

\subsection{Multi-hop QA and supporting-fact selection}

Multi-hop question answering evaluates the ability to combine evidence
across multiple documents. HotpotQA~\cite{yang2018hotpotqa} pairs each
question with ten Wikipedia paragraphs and sentence-level supporting-
fact annotations, making paragraph-selection quality directly
measurable. A broader family of multi-hop benchmarks tests related
phenomena: WikiHop~\cite{welbl2018wikihop} probes cross-document
reasoning chains, 2WikiMultiHopQA~\cite{ho2020constructing}
diagnoses reasoning-step coverage, IIRC~\cite{ferguson2020iirc}
surfaces incomplete-context retrieval, and
MuSiQue~\cite{trivedi2022musique} composes single-hop questions into
provably multi-hop ones to stress-test shortcut learning.
Compositional-gap studies~\cite{press2023measuring} and modular
decompositions~\cite{khot2023decomposed} further isolate where
multi-hop models fail. A line of work frames supporting selection as
a separate stage: cognitive graph retrievers~\cite{ding2019cognitivegraph}
and iterative pathway retrievers~\cite{asai2020learning} traverse
linked entities across paragraphs, while transformer-based sentence
selectors~\cite{nie2019revealing} predict supporting sentences
directly. These methods are accurate but comparatively heavy. Our
TrustScore sits strictly upstream of any such selector: it is
pointwise, scores paragraphs in isolation, and does not reason over
the candidate graph.

\subsection{Retrievers and RAG benchmarks}

Dense passage retrieval~\cite{karpukhin2020dense} replaced classical
sparse retrievers on open-domain QA, and downstream variants including
Fusion-in-Decoder~\cite{izacard2021fid},
ColBERT~\cite{khattab2020colbert,santhanam2022colbertv2}, and
GTR-style generalisable dual encoders~\cite{ni2022gtr} pushed the
retrieval-side recall ceiling further. More recent contrastive
sentence-embedding recipes~\cite{wang2022text} underpin many
zero-shot IR systems. Retrieval-quality benchmarking at scale is
served by BEIR~\cite{thakur2021beir} for zero-shot IR and by the
RGB benchmark~\cite{chen2024rgb} for RAG under noisy, counterfactual,
and conflicting evidence. Classical fact-verification datasets
FEVER~\cite{thorne2018fever} and
FEVEROUS~\cite{aly2021feverous} motivated many of the
evidence-selection metrics we borrow here.

\subsection{Retrieval-augmented generation and context scoring}

Retrieval-augmented generation~\cite{lewis2020rag,guu2020realm} retrieves
a small context set at inference time to ground a
generator, and a growing family of systems explores different
retrieval-generation couplings: RETRO~\cite{borgeaud2022retro} attends
to retrieved chunks during pre-training, Atlas~\cite{izacard2023atlas}
jointly fine-tunes retriever and generator, Active-RAG~\cite{jiang2023active}
retrieves selectively during decoding, REPLUG~\cite{shi2023replug}
treats the generator as a black box, Chain-of-Note~\cite{yu2023chainofnote}
injects per-passage notes to improve robustness, and
RECOMP~\cite{xu2024recomp} compresses retrieved context for
bandwidth-constrained generators. Query-expansion tricks from large
language models~\cite{ma2023queryexpansion} shift some of this burden
upstream. Recent work has probed the failure modes of single-signal
retrieval: LLMs are sensitive to the position of the relevant
context~\cite{liu2024lostinthemiddle}, are easily distracted by
irrelevant context~\cite{shi2023distracted}, struggle with on-topic
distractors~\cite{yoran2024makingrag}, require explicit
context-faithful prompting to avoid prior-knowledge
shortcuts~\cite{zhou2023context}, and degrade further under
conflicting passages~\cite{xie2024adaptive}. Our work complements this
line: rather than train a retriever end-to-end, we ask whether a
\emph{post-retrieval} trust-aware scorer, built from cheap signals and
a small learned combiner, can separate gold from distractor candidates
more reliably than a single relevance signal.

\subsection{Source attribution and trust in LLM outputs}

A parallel thread studies how language-model outputs should credit the
evidence they stand on. GopherCite~\cite{menick2022gophercite} and
Attributed QA~\cite{bohnet2022attributed} fine-tune generators to
produce quotation-backed answers; Gao~et~al.~\cite{gao2023enabling}
survey attribution in generative systems and their
RARR~\cite{gao2023rarr} post-hoc revises model outputs against
retrieved evidence; Rashkin~et~al.~\cite{rashkin2023measuring} define
measurable attribution criteria; factual-consistency evaluation is
re-examined in TRUE~\cite{honovich2022true}; expert-curated
attribution benchmarks such as ExpertQA~\cite{malaviya2023expertqa}
stress-test citation quality in high-stakes domains; and Asai~et~al.'s
Self-RAG~\cite{asai2024selfrag} inserts explicit critique tokens that
assess evidence quality at decode time. Mallen~et~al.~\cite{mallen2023when}
argue that parametric and non-parametric memories deserve different
trust weights depending on question popularity, a framing close in
spirit to ours. These methods operate during or after generation and
presume that a useful evidence pool has already been selected. Our
setting is upstream: we ask which paragraphs deserve to enter the
evidence pool in the first place, and which signals predict that.

\subsection{Confidence, consistency, and hallucination filtering}

Trust-aware selection is closely related to confidence estimation and
hallucination detection. Self-consistency~\cite{wang2023selfconsistency}
uses agreement across sampled reasoning chains as a correctness
signal; selfcheckGPT~\cite{manakul2023selfcheckgpt} extends this to
zero-resource hallucination detection;
Kadavath~et~al.~\cite{kadavath2022language} probe language models'
internal calibration; Jiang~et~al.~\cite{jiang2021know} study
calibration specifically for QA; and
Tian~et~al.~\cite{tian2023justask} show that simply asking for a
calibrated confidence outperforms model log-probabilities after RLHF.
Hallucination is an active survey area in its own
right~\cite{ji2023survey,zhang2023siren,li2024evaluating}. These
methods rely on consistency across model samples. Our consistency
signal is structurally similar but operates across candidate
\emph{passages} at retrieval time, and our learned combiner discovers
that in a distractor-rich pool this signal is anti-correlated with
gold, reversing the naive reading of consistency as corroboration.

\subsection{Distinguishing contribution}

Most closely related is the broader idea, articulated in recent
position papers~\cite{wu2024retrieval,ram2023incontextrag}, that RAG
pipelines should treat sources as heterogeneous entities with
per-source trust rather than a homogeneous relevance pool. These
discussions are largely conceptual. We operationalize the idea on a
standard benchmark and show two concrete results. First, naive uniform
trust-signal averaging is \emph{actively harmful} on HotpotQA
distractor: not every trust heuristic is positively correlated with
gold, and averaging them without calibration is worse than using
relevance alone. Second, a six-feature logistic regression fit on 100
questions per seed recovers the correct sign of each signal, beats
SBERT by 3.3 absolute Precision@2 points, and more than doubles the
gold-versus-distractor score gap. The ablation isolates a concrete
instance of this phenomenon in HotpotQA distractor, and we are not
aware of an existing empirical demonstration at this granularity.
